\section{Environmental Triggers of Bias in RAG}
\label{sec:triggers}
% Length: ~20%
% Focus: Analyze the specific external pressures that cause the judge to fail. Do not list papers sequentially; group them by trigger type.

% Must contain:
% 1. Context Length: How massive input sizes destroy evaluation accuracy.
% 2. Distractors and Noise: How retrieved irrelevant information confuses the judge.
While standard LLM evaluators exhibit well-documented vulnerabilities such as verbosity and self-enhancement biases \citep{wang2023fair}, RAG frameworks introduce additional external challenges. These environmental factors compromise the judge's evaluation accuracy. We identify two primary triggers that systematically degrade the reliability of LLM-as-a-Judge in RAG systems: extreme context length and unavoidable retrieval noise.

\subsection{Context Length and Positional Bias}

In the RAG evaluation workflow, a judge must process multiple lengthy retrieved documents simultaneously, often approaching the limits of its context window. Research has extensively documented the ``Lost in the Middle'' phenomenon, where language models exhibit a U-shaped performance curve. They reliably process information at the beginning or end of a prompt but fail to utilize content in the middle \citep{lostinthemiddle2023}. When applied to automated evaluation, this positional vulnerability translates into scoring biases. The physical distance between related pieces of information within a prompt can measurably affect the judge's assessment \citep{distance2024}. In RAG-QA settings, studies show that evaluators often assign higher scores to responses relying on the top-ranked retrieved document, overlooking valid or contradictory evidence buried deeper in the context \citep{ragqa2024, posbias2025}. Although structural interventions such as calibrating positional attention \citep{foundinmiddle2024} provide partial remedies, performance degradation under extreme context lengths remains a primary cause of evaluation failure.

\subsection{Distractors, Noise, and Cognitive Overload}

Beyond context length, the quality of retrieved content introduces the second major environmental trigger: noise. Because retrieval algorithms like dense vector search are fundamentally probabilistic, the context presented to the judge is rarely clean. It typically contains contextual distractors, which are pieces of information semantically similar to the query but factually unrelated.

The presence of these distractors severely impairs the reasoning capabilities of the judge. When forced to navigate such noise, models experience systematic performance degradation, a phenomenon termed ``Lost in the Noise'' \citep{lostinnoise2026}. Critically, the evaluator faces a dual burden. It must parse the noisy retrieved context and simultaneously use that same context to verify the faithfulness of the generated response. Analysis of internal representations confirms that distracting tokens can disrupt the model's reasoning pathways \citep{internalrep2026}. This disruption often causes the judge to penalize a correct answer due to the presence of surrounding irrelevant information.

To counter this issue, researchers have explored dynamic context selection methods that filter distractors before they reach the evaluator \citep{dynamiccontext2025}. However, when noise does enter the prompt, it forces the model to choose between trusting the flawed external context and relying on its pre-trained memory. This specific tension gives rise to the knowledge conflicts examined in the following section.